\section{Architecture Characteristics}\label{sec:3}

\subsection{Identifying Candidate Architecture Characteristics}\label{sec:3.1}

Architecture characteristics express system success criteria.
The course requires a characteristic to concern quality beyond domain
functionality, influence design structure and matter to system success. Candidates were extracted from explicit requirements,
domain concerns and implicit knowledge. This
distinguishes qualities such as Security from tactics such as RBAC, and
Reliability from the choice of broker/database.

{\small
\begin{longtable}[]{|>{\raggedright\arraybackslash}p{\dimexpr 0.3234\linewidth-2\tabcolsep\relax}|>{\raggedright\arraybackslash}p{\dimexpr 0.3233\linewidth-2\tabcolsep\relax}|>{\raggedright\arraybackslash}p{\dimexpr 0.3233\linewidth-2\tabcolsep\relax}|}
\hline
\textbf{Candidate} & \textbf{Requirement origin} & \textbf{Structural influence} \\ \hline
\endhead
AC-01 Performance & E per-read deadline and workload & Completion
placement, queue scheduling and resource isolation \\ \hline
AC-02 Reliability/durability & E movement/health no-loss; P fault
extension & Persistence, commit/ACK boundaries, idempotency and
correction history \\ \hline
AC-03 Availability/continuity & I outage risk; P offline scope & Local
processing, cached context, pending storage and recovery \\ \hline
AC-04 Security & E RBAC; P scopes/guards & IAM, resource authorization
and human/machine identities \\ \hline
AC-05 Auditability & E comprehensive audit; P evidence policy & Source
evidence, collection, protection and traceable access \\ \hline
AC-06 Interoperability & E API exchange; I/P quality & Adapters, schema
validation and durable submission lifecycle \\ \hline
AC-07 Scalability/configurability & E current scale; I/P expansion &
Resource capacity, namespaces, configuration and provisioning \\ \hline
AC-08 Maintainability/testability & E SOLID and M2; I/P change criteria
& Stable core interfaces, replaceable adapters and deterministic
tests \\ \hline
\end{longtable}
}

Availability is a Reliability sub-characteristic in the course taxonomy. It has a separate criterion here so continuity is
measured independently of record preservation. Auditability elaborates
Accountability and is measured separately from permission correctness.
Semantic integrity and operational observability support these
candidates through correct identity/count interpretation and diagnosable
processing.

\subsection{Success Criteria}\label{sec:3.2}

Each criterion combines its characteristic, source, measurable outcome
and planned verification. The M1 criteria remain design contracts.
\Cref{sec:m2-health-results} records the subset exercised in the M2
Health backend.

{\small
\begin{longtable}[]{|>{\raggedright\arraybackslash}p{\dimexpr 0.2426\linewidth-2\tabcolsep\relax}|>{\raggedright\arraybackslash}p{\dimexpr 0.2425\linewidth-2\tabcolsep\relax}|>{\raggedright\arraybackslash}p{\dimexpr 0.2425\linewidth-2\tabcolsep\relax}|>{\raggedright\arraybackslash}p{\dimexpr 0.2425\linewidth-2\tabcolsep\relax}|}
\hline
\textbf{Criterion} & \textbf{Characteristic and source} & \textbf{Measurable acceptance} & \textbf{Planned verification} \\ \hline
\endhead
SC-01 & Performance; E/P, NFR-01 & Every required read completes
correctly within \(\leq\)2 seconds; zero missing outcomes or deadline
violations & Full input/outcome reconciliation and source/completion
timing, including uncertainty and failures \\ \hline
SC-02 & Reliability; E/P, NFR-02 & Zero missing committed
movement/health IDs or contents; no duplicate business effects within
the declared envelope & Input/history comparison; uncertain-response and
crash/replay controls \\ \hline
SC-03 & Availability; I/P, NFR-03 & Local gate core survives the
60-minute test; its backlog completes within 600 seconds of
recovery-ready while live traffic continues & Outage/live cohort IDs,
local timing, backlog and final central effects/evidence \\ \hline
SC-04 & Security; E/P, NFR-04 & All expected-deny cases reject without
protected disclosure or unauthorized effects; authorized controls
succeed & Role/action matrix plus wrong-scope, revocation, transfer and
export checks \\ \hline
SC-05 & Auditability; E/P, NFR-05 & Valid durable source evidence for
100\% of covered committed operations; required access events covered;
complete canonical set after recovery & Expected ID/content manifests;
source/collector faults and protected access controls \\ \hline
SC-06 & Interoperability; E/I/P, NFR-06 & Zero lost persisted intents;
valid approved payloads; no unsafe retries; completion requires
contractual acceptance evidence & API mocks for final success, 202, rate
limits and timeout-after-effect \\ \hline
SC-07 & Scalability; E/I/P, NFR-07 & Current dataset and declared
workload supported; compatible additions need zero core-rule edits and
preserve scopes/data & Dataset/load checks and preapproved
configuration-change demonstration \\ \hline
SC-08 & Maintainability; I/P, NFR-08 & Two compatible reader/API adapter
changes need zero core-rule edits; relevant contract/regression tests
pass & Frozen core-file diffs, interface review and repeatable domain
tests \\ \hline
\end{longtable}
}

Record reconciliation compares identities and contents, since matching
totals can hide a missing record and a duplicate. Availability recovery
is complete when the declared business/projection/evidence stages finish
the outage cohort; an empty queue alone is insufficient. External mock
results establish only the declared mock contract, not real National API
compatibility.

\subsection{Driving Characteristics}\label{sec:3.3}

The three driving characteristics are D-01 Reliability / Durability, D-02
Performance and D-03 Security / Auditability. Each combines a mandatory
outcome with independent structural consequences, and the course
recommends prioritization and keeping the primary list short. All three
are rated High importance and High difficulty: failing any of them
violates a mandatory requirement, and their interactions shape
persistence, timing and authorization. These ratings are design
assessments, not measured results.

\subsubsection{D01 Reliability and Durability}\label{sec:3.3.1}

Lost or duplicated movements corrupt counts, while lost health records
undermine complete profiles. This driver requires durable ownership,
atomic local effects/evidence, stable operation identities and preserved
correction history. It therefore affects stores, communication contracts
and recovery responsibilities. Additional writes, storage and
reconciliation are the trade-offs. Confirmation and ACK must match the
stage actually completed.

\subsubsection{D02 Performance}\label{sec:3.3.2}

The two-second deadline influences where gate processing completes and
which dependencies lie on its critical path. A low average passage rate
cannot rule out burst queues. Performance therefore remains a primary
driver even when average throughput seems modest.

The selected design completes the gate outcome locally and handles
central synchronization, reports and external delivery separately.
CPU/storage contention still requires resource control; local placement
alone proves no timing guarantee. The trade-off is balancing durable
work, concurrency and operating resources while retaining correctness
and audit. This requirement does not independently mandate
microservices.

\subsubsection{D03 Security and Auditability}\label{sec:3.3.3}

RBAC and comprehensive audit affect all protected operations. This
driver determines identity/resource boundaries, authorization placement,
workflow approvals and evidence ownership. Its costs include
authorization and durable-evidence work and restricted offline human
operations. Permission correctness and audit completeness remain
separate measures even though they share a driver group.

Availability remains significant for the selected edge/offline design,
with its own NFR, criterion and scenario. It is outside the primary
three because the assignment gives a mandatory timing constraint but no
outage duration or production availability target. Distributed styles
introduce latency, consistency and failure trade-offs, so \Cref{sec:4}
compares styles through concrete variants against these three drivers.

\subsection{Quality Attribute Scenarios}\label{sec:3.4}

Each scenario contains Source, Stimulus, Environment, Artifact, Response
and Response Measure, following the SEI formulation. QAS01--08 retain
the M1 scenarios. QAS09 and QAS10 extend them for on-site mobile
recording, where intermittent connectivity and cattle identification
errors can interrupt a clinical workflow.\cite{assignmentM2} Fixture
budgets are project-defined (P). \Cref{tab:m2-health-cases} records the
specific backend cases exercised in M2.

\subsubsection{QAS01 Correct and Timely Gate Processing}\label{sec:3.4.1}

{\small
\begin{longtable}[]{|>{\raggedright\arraybackslash}p{\dimexpr 0.4849\linewidth-2\tabcolsep\relax}|>{\raggedright\arraybackslash}p{\dimexpr 0.4851\linewidth-2\tabcolsep\relax}|}
\hline
\textbf{Element} & \textbf{Specification} \\ \hline
\endhead
Source & Qualified RFID reader or emulator with independently known
inputs, direction and source timing \\ \hline
Stimulus & Approximately 200 cattle pass one farm' s
gates in 30 minutes; declare raw-read multiplicity and burst schedule \\ \hline
Environment & Normal operation, healthy edge/storage, valid context and
current-scale dataset; resources declared \\ \hline
Artifact & Reader adapter and local movement/count/evidence
processing \\ \hline
Response & Commit each correct logical movement once; retain valid
duplicate/invalid-input dispositions and corresponding count effects \\ \hline
Response Measure & Every required read completes within \(\leq\)2
seconds from source creation, including clock uncertainty; zero
missing/incorrect outcomes or duplicate effects \\ \hline
\end{longtable}
}

Trace: NFR-01/02/05; SC-01/02/05; D-02 and D-01. Repeated observations
do not automatically represent additional cattle passages.

\subsubsection{QAS02 WAN Outage and Recovery}\label{sec:3.4.2}

{\small
\begin{longtable}[]{|>{\raggedright\arraybackslash}p{\dimexpr 0.4849\linewidth-2\tabcolsep\relax}|>{\raggedright\arraybackslash}p{\dimexpr 0.4851\linewidth-2\tabcolsep\relax}|}
\hline
\textbf{Element} & \textbf{Specification} \\ \hline
\endhead
Source & Network fault controller with provisioned reader/worker
identities \\ \hline
Stimulus & Disconnect central connectivity for 60 minutes; generate 400
valid movements/farm and restore services while live traffic
continues \\ \hline
Environment & Test with healthy local power/storage, sufficient
buffer and valid cached context \\ \hline
Artifact & Edge gate core/pending storage, Sync and central
business/count/evidence stages \\ \hline
Response & Continue local processing, identify stale context, retain
pending work and replay safely without blocking live gate processing \\ \hline
Response Measure & Preserve all 400 movements/farm without duplicate
effects; local reads retain SC-01; outage cohort finishes within 600
seconds of actual recovery-ready \\ \hline
\end{longtable}
}

Trace: NFR-03/01/02; SC-03/01/02; AC-03. Offline human operations follow
\Cref{sec:2.2.2}; three-farm testing uses 1,200 outage movements. The
budget applies to this finite cohort.

\subsubsection{QAS03 Consumer Crash and Replay}\label{sec:3.4.3}

{\small
\begin{longtable}[]{|>{\raggedright\arraybackslash}p{\dimexpr 0.4849\linewidth-2\tabcolsep\relax}|>{\raggedright\arraybackslash}p{\dimexpr 0.4851\linewidth-2\tabcolsep\relax}|}
\hline
\textbf{Element} & \textbf{Specification} \\ \hline
\endhead
Source & Fault controller and duplicate publisher \\ \hline
Stimulus & Crash after store commit but before ACK, restart and replay
the same 100 distinct normalized movement IDs \\ \hline
Environment & Intact storage; 100 distinct registered cattle initially
outside the barn; initial count zero; no new input. Repeat for
separately stored projections. \\ \hline
Artifact & Movement/count transactions, idempotency markers and required
source evidence \\ \hline
Response & Return the existing committed outcome; each store protects
its own marker/effect/evidence and acknowledges after commit \\ \hline
Response Measure & Exactly 100 IN movement/count effects and final count
100; zero missing/extra effects or required canonical facts within a
60 seconds of actual consumer recovery-ready \\ \hline
\end{longtable}
}

Trace: NFR-02/05; SC-02/05; D-01. Before-commit failure rolls back
together; same ID with different content is a conflict. The 60-second
budget is for this finite test cohort, not a production RTO or global
transaction guarantee.

\subsubsection{QAS04 Unauthorized Role or Scope}\label{sec:3.4.4}

{\small
\begin{longtable}[]{|>{\raggedright\arraybackslash}p{\dimexpr 0.4849\linewidth-2\tabcolsep\relax}|>{\raggedright\arraybackslash}p{\dimexpr 0.4851\linewidth-2\tabcolsep\relax}|}
\hline
\textbf{Element} & \textbf{Specification} \\ \hline
\endhead
Source & Authenticated BS/FM/SA/EM with a declared role and scope \\ \hline
Stimulus & Request protected data or mutation outside permission:
wrong-farm access, SA clinical write and revoked export. During WAN loss
also attempt online-only writes/new login and local dashboard reads
before and after lease expiry. \\ \hline
Environment & Online: authoritative policy and known allow/deny
fixtures. Offline: declared cached grant/scope and a known policy-lease
expiry; central revocation is observable only after reconnection or
lease expiry. \\ \hline
Artifact & Resource services, history queries, report execution/download
and audit source \\ \hline
Response & Enforce action/resource scope at the owner. Online checks use
current policy. Offline, only existing leased read-only dashboard access
and provisioned machine capture remain eligible; deny online-only
operations and expired leases. Retain required security evidence. \\ \hline
Response Measure & All declared allow/deny decisions are correct; denied
cases disclose no protected data and cause no unauthorized effect.
In-scope local reads succeed before the permitted expiry and are denied
at or after it. Online-only operations/new login fail during WAN loss;
audit gaps remain explicit. \\ \hline
\end{longtable}
}

Trace: NFR-04/05; SC-04/05; D-03; ADR-07/12. Test role/scope and
online/offline boundaries; gateway checks alone are insufficient.
Offline revocation exposure is bounded by the declared read-only lease,
rather than assumed instantaneous.

\subsubsection{QAS05 Health Write and Audit Failure}\label{sec:3.4.5}

{\small
\begin{longtable}[]{|>{\raggedright\arraybackslash}p{\dimexpr 0.4849\linewidth-2\tabcolsep\relax}|>{\raggedright\arraybackslash}p{\dimexpr 0.4851\linewidth-2\tabcolsep\relax}|}
\hline
\textbf{Element} & \textbf{Specification} \\ \hline
\endhead
Source & Authorized BS/FM and a fault controller \\ \hline
Stimulus & Submit one health intent and three same-payload retries while
either the remote collector is unavailable or required source evidence
insertion fails \\ \hline
Environment & Online health workflow; case A has healthy source storage,
case B fails the mandatory source evidence step \\ \hline
Artifact & Health version, operation receipt, source evidence/outbox and
collector \\ \hline
Response & A commits once with source evidence and relays later; B rolls
back and reports no success; resolve lost responses with the same
authorized intent \\ \hline
Response Measure & One intended health effect after a successful commit,
zero partial committed effects in B before repair; complete history and
source evidence; the small canonical audit cohort completes within a
30 seconds of collector-ready \\ \hline
\end{longtable}
}

Trace: NFR-02/05/04; SC-02/05/04; D-01/D-03. Reauthorization and
matching prior receipts prevent a retry from becoming a second
treatment. The 30-second budget applies to this finite cohort, not a
global audit-delivery SLA.

\subsubsection{QAS06 Ambiguous External Delivery}\label{sec:3.4.6}

{\small
\begin{longtable}[]{|>{\raggedright\arraybackslash}p{\dimexpr 0.4849\linewidth-2\tabcolsep\relax}|>{\raggedright\arraybackslash}p{\dimexpr 0.4851\linewidth-2\tabcolsep\relax}|}
\hline
\textbf{Element} & \textbf{Specification} \\ \hline
\endhead
Source & External API mock and network fault controller \\ \hline
Stimulus & Return 202, rate limits or timeout after a receiver side
effect; provide separate contracts with and without idempotency \\ \hline
Environment & Persisted approved intent/schema; receiver safety
semantics explicitly declared; gate processing independent \\ \hline
Artifact & Integration adapter, durable intent/attempt states and
reconciliation path \\ \hline
Response & Retain pending/unknown outcomes; respect applicable
Retry-After and bounded retries; retry ambiguous operations only when
contractually safe \\ \hline
Response Measure & Zero lost intents; one receiver effect in the
idempotent fixture; no blind ambiguous retry without safety; Delivered
only after final business acceptance evidence \\ \hline
\end{longtable}
}

Trace: NFR-06/02; SC-06/02; AC-06. Retry-budget exhaustion retains
unresolved work. Mocks cannot establish actual National API
capabilities.

\subsubsection{QAS07 Compatible Configuration Growth}\label{sec:3.4.7}

{\small
\begin{longtable}[]{|>{\raggedright\arraybackslash}p{\dimexpr 0.4849\linewidth-2\tabcolsep\relax}|>{\raggedright\arraybackslash}p{\dimexpr 0.4851\linewidth-2\tabcolsep\relax}|}
\hline
\textbf{Element} & \textbf{Specification} \\ \hline
\endhead
Source & Deployment team using a preapproved farm/barn/reader
provisioning package \\ \hline
Stimulus & Add a compatible farm namespace and reader bindings through
configuration \\ \hline
Environment & Current three-farm dataset; synthetic additional farm,
approved identities/scopes and allocated resources \\ \hline
Artifact & Registry/configuration, routing, scope policy and core
business modules \\ \hline
Response & Validate/provision bindings, preserve existing records and
isolate the new scope \\ \hline
Response Measure & Zero edits to predeclared core movement/count/health
rules; correct positive/negative scope checks; existing workload retains
SC-01 before and after change \\ \hline
\end{longtable}
}

Trace: NFR-07/01/04; SC-07/01/04; AC-07. This verifies compatible
onboarding, not unlimited capacity or an undefined runtime creation
permission.

\subsubsection{QAS08 Compatible Adapter Changes}\label{sec:3.4.8}

{\small
\begin{longtable}[]{|>{\raggedright\arraybackslash}p{\dimexpr 0.4849\linewidth-2\tabcolsep\relax}|>{\raggedright\arraybackslash}p{\dimexpr 0.4851\linewidth-2\tabcolsep\relax}|}
\hline
\textbf{Element} & \textbf{Specification} \\ \hline
\endhead
Source & Developer in development/CI \\ \hline
Stimulus & Independently replace a reader protocol adapter and an API
adapter while preserving their domain-facing contracts \\ \hline
Environment & M2 implementation, frozen core-file inventory and
deterministic domain/contract fixtures \\ \hline
Artifact & Adapter interfaces/configuration and core
movement/count/health modules \\ \hline
Response & Implement changes at the adapter boundary and run relevant
contract/domain/regression tests \\ \hline
Response Measure & Both compatible changes require zero core-rule edits;
relevant tests pass and repeated deterministic fixtures yield expected
outcomes \\ \hline
\end{longtable}
}

Trace: NFR-08; SC-08; AC-08 and SOLID. Domain-semantic changes need
separate analysis; absent implementation produces Not Run.

\newpage
\subsubsection{QAS09 Mobile Capture and Restart}\label{sec:qas09}

{\small
\begin{longtable}{|p{0.19\linewidth}|p{0.72\linewidth}|}
\hline
\textbf{Element} & \textbf{Specification} \\ \hline
\endhead
Source & An authorized barn supervisor uses a mobile client and
reader. A tester interrupts connectivity or restarts the app. \\ \hline
Stimulus & Scan and confirm the animal, enter the clinical details, and
attempt to save. Interrupt connectivity before or after saving, then
reconnect and retry with the same intent ID. \\ \hline
Environment & The tag and custody start in the verified scope. The
client loses and regains network access. Option A follows the online-only
baseline. Option B proposes durable local drafts during the outage. \\ \hline
Artifact & Mobile form, local draft store and sync state, together with
the Health intent, receipt, clinical record and source evidence. \\ \hline
Response & Under A, the app reports that a new offline save is unavailable.
Under B, the app reports a saved draft only after the local commit and
restores the same cattle, actor, intent and content after restart. Upon
reconnection, Health checks current authorization, custody and version.
The app shows acceptance only after the server confirms the result. \\ \hline
Response\newline Measure & Zero server acceptances during the outage and zero
false saved or accepted states. For B, the proposed fixture restores
20 of 20 acknowledged drafts. With unchanged rights and reconnection
ready, the valid cohort yields 20 matching effects, receipts and evidence
rows within 120 seconds (P), with zero duplicates or wrong-cattle effects.
Conflicts retain their actual nonaccepted outcome. \\ \hline
\end{longtable}
}

Trace: SC-02 and SC-05, \Cref{adr:12}, and T05/T06. Offline drafts and
their proposed measures remain design-only pending confirmation of
scope. The online retry controls are recorded in
\Cref{tab:m2-health-cases}.

\newpage
\subsubsection{QAS10 Cattle Identity and Current Authorization}\label{sec:qas10}

{\small
\begin{longtable}{|p{0.19\linewidth}|p{0.72\linewidth}|}
\hline
\textbf{Element} & \textbf{Specification} \\ \hline
\endhead
Source & A barn supervisor or farm manager, an out-of-scope user,
the RFID reader and an administrator who can change grants or custody. \\ \hline
Stimulus & Scan an unknown, deactivated or ambiguous tag. Attempt a
submission, retry, or receipt or history read after a grant or custody
change, or submit with a stale resource version. \\ \hline
Environment & Health is online with server-owned cattle and custody data.
Allow and deny cases declare the caller's current grants, resource
version and expected cattle binding. \\ \hline
Artifact & Scan and form selection, verified principal and grants,
canonical cattle and record context, Health policy and clinical outcome. \\ \hline
Response & Require explicit cattle confirmation before capture. An unknown
tag stays unresolved. A known deactivated tag still resolves to cattle
identity and history for authorized callers, and cannot start a new
clinical write. Each protected operation checks current rights. A
new submission with a stale version goes to review when scope still
allows it. A committed retry follows current receipt authorization.
An out-of-scope lookup gives no cattle data. \\ \hline
Response\newline Measure & Zero unauthorized effects or protected-data
disclosures in deny cases. Unknown tags create zero cattle records.
For allowed flows, the confirmed cattle and submitted content match
the resulting receipt and history. A stale uncommitted intent produces
zero clinical records, and an authorized replay produces no extra
effect. \\ \hline
\end{longtable}
}

Trace: SC-02, SC-04 and SC-05, \Cref{adr:07,adr:08}, and T06--T09.
\Cref{tab:m2-health-cases} covers the backend examples exercised on
\texttt{staging}. The mobile confirmation flow and integration with
the real identity and grant providers remain pending.
